% Resultados -- RQ 07 e RQ 08 (parte do Pedro).
% Números de resultados/rq7_*.csv e resultados/rq8_*.csv.

\subsection{RQ 07: sensibilidade à definição de deploy}

Recalculamos a classificação DORA dos 350 repositórios sob três operacionalizações de
\textit{deploy}: \textbf{C1} (referência), que usa apenas \textit{releases}, com
\textit{lead time} (a) e CFR (a); \textbf{C2}, que inclui pré-\textit{releases} e adota as
variantes (b); e \textbf{C3}, que usa \textit{tags} com o CFR na variante (a).

\begin{table}[ht]
\centering
\caption{Distribuição das categorias DORA sob as três definições de \textit{deploy}.}
\label{tab:rq7dist}
\begin{tabular}{lrrr}
\hline
\textbf{Categoria} & \textbf{C1} & \textbf{C2} & \textbf{C3} \\
\hline
Elite  &  25 &   8 &  33 \\
High   & 143 & 134 & 164 \\
Medium & 159 & 186 & 138 \\
Low    &  23 &  22 &  15 \\
\hline
\end{tabular}
\end{table}

A concordância entre as classificações, medida pelo kappa ponderado, fica na faixa moderada
(Tabela~\ref{tab:rq7pares}). Entre um quarto e um terço dos repositórios trocam de categoria
conforme a definição adotada. A faixa \textit{Elite} é a mais sensível: vai de 8 repositórios
em C2 a 33 em C3, uma variação de quatro vezes sobre o mesmo conjunto de projetos e a mesma
janela de observação.

\begin{table}[ht]
\centering
\caption{Concordância entre as classificações, par a par.}
\label{tab:rq7pares}
\begin{tabular}{lrr}
\hline
\textbf{Par} & \textbf{Repositórios que mudam} & \textbf{Kappa ponderado} \\
\hline
C1 $\times$ C2 & 32,9\% & 0,515 \\
C1 $\times$ C3 & 25,4\% & 0,657 \\
C2 $\times$ C3 & 28,0\% & 0,602 \\
\hline
\end{tabular}
\end{table}

\subsection{RQ 08: rework rate}

O \textit{rework rate} é a proporção de \textit{releases} corretivas publicadas na janela,
calculada como \texttt{cfrb\_releases\_corretivas} dividido por \texttt{n\_releases}.
Diferente do CFR (b), ele conta toda \textit{release} corretiva, e não apenas as que vieram
em até 7 dias depois de outra.

\begin{table}[ht]
\centering
\caption{\textit{Rework rate} e CFR (b) nos 350 repositórios.}
\label{tab:rq8}
\begin{tabular}{lrrrr}
\hline
\textbf{Métrica} & \textbf{n} & \textbf{Q1} & \textbf{Mediana} & \textbf{Q3} \\
\hline
\textit{Rework rate} & 350 & 0,225 & 0,500 & 0,726 \\
CFR (b)              & 350 & 0,052 & 0,200 & 0,416 \\
\hline
\end{tabular}
\end{table}

A mediana do \textit{rework rate} é 2,5 vezes a do CFR (b). As duas medidas se correlacionam
($\rho = 0{,}522$, $p < 0{,}001$), mas \textbf{309 dos 350 repositórios têm \textit{rework
rate} maior que o CFR (b)}, e 21 não publicaram nenhuma \textit{release} corretiva na janela.

A leitura é direta: metade das \textit{releases} desses projetos corrige alguma coisa, mas
apenas uma em cada cinco corrige algo publicado na semana anterior. A maior parte do
retrabalho é manutenção de rotina, e não resposta a uma entrega recém-quebrada --- o que
reforça que o CFR (b), ao exigir a janela de 7 dias, mede um fenômeno mais estreito que
``\textit{release} corretiva''.
